
The runaway loss rate $\overline{\Gamma }_{R}$ is also an important physical
quantity that is determined by the code. It becomes significant when the
Ohmic electric field is large. For the benchmarking procedure, the Ohmic
electric field is then set to $\overline{E}=0.08.$ All other parameters are
kept constant as compared to the benchmarking of the Ohmic conductivity done
in Sec. \ref{conducsection}. The corresponding tokamak parameter file is
nevertheless named ``CQL3D\_RUNAWAY'', since it corresponds to conditions
that have been used for code comparison with the well know CQL3D program (%
\cite{har93c1_IAEA}). By definition, the code is no more conservative since
runaway electrons are leaving the domain of integration. Consequently, the
particle loss must be compensated in order to keep the totalnumber of
particle constant. In addition, in $3-D$ mode, there are numerical losses,
as shown in the previous section, that overcome physical ones. Consequently,
the same procedure is used and the Maxwellian solution is enforced at $i=1/2$
while the density is normalized at each iteration. It has been cros-checked
that the compensation of losses has no influence on the final solution and
therefore the runaway rate.

In a first step, a local analysis is performed, without considering bounce
averaging. Here, only the Maxwellian electron-electron collision operator is
considered, for allowing comparisons with analytical solutions, as shown in 
\ref{table_run1}

\begin{table}[tbp]
\begin{center}
\begin{tabular}{|l|l|l|}
\hline
$Z_{eff}$ & $\overline{\Gamma }_{R}$ [DKE code] & $\overline{\Gamma }_{R}$
[Kulsrud et al.] \\ \hline
$1$ & $2.8487\times 10^{-4}$ & $3.177\times 10^{-4}$ \\ \hline
$2$ & $1.5904\times 10^{-4}$ & $1.735\times 10^{-4}$ \\ \hline
$3$ & $9.6565\times 10^{-5}$ & $1.047\times 10^{-4}$ \\ \hline
$10$ & $8.1878\times 10^{-6}$ & $9.0\times 10^{-6}$ \\ \hline
\end{tabular}
\end{center}
\caption{Runaway rate as function of the effective charge using the
Maxwellian e-e collision model}
\label{table_run1}
\end{table}

The agreement between numerical results obtained with the drift kinetic code
and the code written by Kulsrud and coauthors is reasonably good \cite{kul73}%
, a systematic difference of $-15\%$ being observed  for all $Z_{eff}$
values. For $Z_{eff}=1$, the agreement is slightly better with the model of
Kruskal-Bernstein \cite{kul73}, and the relative difference is now positive
and less than $7\%$. The role played by the numerical grid is marginal, less
than $1\%$. when  $n_{p}$ is varying from $88\ $to $125\ $or $165$ while $%
n_{\xi _{0}}$ is kept constant at $120.$

A full $3-D$ calculation has been performed with flat profiles at normalized
radial positions on the spatial flux grid $\left[ 0.2,0.4,0.6,0.8\right] ,$
giving the following radial positions $\rho =\left[
0.14142,0.31623,0.5099,0.70711,0.90554\right] $. In that case, the code
calculates itself the pitch-angle grid and the number of step is $n_{\xi
_{0}}=168,$ while $n_{p}=125.$The runaway rate is found to be independent of
the method of calculations, and the Kruskal-Bernstein is well recovered when 
$Z_{eff}=1$ at all plasma radii. 

Bounce averaged calculations have be performed first locally. As shown in
Fig. \ref{ohmrun}, it is observed $\overline{\Gamma }_{R}$ decreases with
the inverse aspect ratio $\epsilon .$ The influence of toroidicity on
Dreicer generation, which is not predicted by the theory of Gurevitch \cite
{gur61}, is in good agreement with never published calculations using the
CQL3D code, but also simulations of runaways during tokamak disruptions \cite
{eri03}. It is shown in Ref. \cite{eri03} that trapping effects may be
approximately described by the relation
\begin{equation}
\overline{\Gamma }_{R}\left( \epsilon \right) \simeq \overline{\Gamma }%
_{R}\left( \epsilon =0\right) \sqrt{1-2\epsilon }
\end{equation}
when $\epsilon \leq 0.5\ $while  $\overline{\Gamma }_{R}\left( \epsilon \geq %
0.5\right) =0.$ The drift kinetic code confirms therefore this important
result, that the runaway generation is nearly cancelled when $\epsilon \geq %
0.5.$ 

Similar results are obtained in $3-D$ mode, whatever the normalization
reference of the electric field, at the plasma center, or at the local
position. As expected, the runaway rate $\overline{\Gamma }_{R}$ is nearly
independent of the sign of the electric field, and moreover, but also of the
method (analytical or numerical) for calculating the bounce integrals. The
relative accuracy which is less than 1\% close to the plasma center, tends
to decrease down to 50\% at the plasma edge, but since $\overline{\Gamma }%
_{R}$ is very small the error has only a weak influence\footnote{%
It is important to note that the order of magnitude is always correct.}. 

%TCIMACRO{
%\TeXButton{Fig_runaway}{\fig{figures/runaway.eps}{Normalized Ohmic runaway rate as function of the inverse aspect ratio $\epsilon$}{ohmrun}{7cm} }}
%BeginExpansion
\fig{figures/runaway.eps}{Normalized Ohmic runaway rate as function of the inverse aspect ratio $\epsilon$}{ohmrun}{7cm} %
%EndExpansion

Once more, this demonstrates that numerical calculations of the bounce
integrals is very accurate, even for finite inverse aspect ratio. The drop
tolerance has been lowered in that case to $10^{-4}$ for ensuring a stable
convergence, leading to a modest increase of the matrix sizes. However the
rate of convergence is not affected, and the final result is obtained
usually after $11-13$ iterations, even when several radial positions are
considered.
